% GENERATED FILE. Edit canonical lessons, then run npm run generate:books.
% canonical-source-hash: fnv1a64:f32ab56667966c82
% canonical-lessons: RU-C175-tyanut, RU-C175-podnimat, RU-C175-vklyuchat, RU-C175-vyklyuchat, RU-C175-zazhigat

\chapter{To pull, To lift, To switch on, To switch off, To light}
\label{ch:ru-to-pull-to-lift-to-switch-on-to-switch-off-to-light}
\hlchaptermodality{175}

\begin{chapteropening}
\textbf{By the end of this chapter:} \emph{I can say to pull, to lift, to switch on, to switch off and to light.}

Complete the last lesson of chapter 175: I can say to pull, to lift, to switch on, to switch off and to light.
\end{chapteropening}

\section[\texorpdfstring{\ru{тянуть} (tyanút)}{тянуть (tyanút)}]{\ru{тянуть} (tyanút) — to pull}

\label{lesson:RU-C175-tyanut}



\begin{quote}
Before the new one: say the Russian for to translate, then the Russian for to remember.
\end{quote}

\subsection*{You'll want to know: \ru{тянуть}}
\textbf{\ru{тянуть}} (\emph{tyanút}) — \textquotedblleft{}to pull\textquotedblright{}.

\begin{grammarlens}[title={What you've built}]
One more word for this chapter.
\end{grammarlens}

\subsection*{Your turn}
\begin{itemize}
\raggedright
  \item \emph{Say it:} \emph{tyanút}
  \item \emph{Say it:} \emph{tyanút}, once more
  \item \emph{Say it:} say \emph{tyanút}
  \item \emph{Recall:} say the Russian for to translate, then the Russian for to remember, then say \emph{tyanút} again
\end{itemize}

\begin{tcolorbox}[breakable,colback=teal!4,colframe=teal!35!black,title={Before you move on}]
What does \ru{тянуть} mean? (\textquotedblleft{}To pull\textquotedblright{}.) Say it once more.
\end{tcolorbox}
\section[\texorpdfstring{\ru{поднимать} (podnimát)}{поднимать (podnimát)}]{\ru{поднимать} (podnimát) — to lift, to raise}

\label{lesson:RU-C175-podnimat}



\begin{quote}
Before the new one: say the Russian for to remember, then the Russian for to pull.
\end{quote}

\subsection*{You'll want to know: \ru{поднимать}}
\textbf{\ru{поднимать}} (\emph{podnimát}) — \textquotedblleft{}to lift, to raise\textquotedblright{}.

\begin{grammarlens}[title={What you've built}]
One more word for this chapter.
\end{grammarlens}

\subsection*{Your turn}
\begin{itemize}
\raggedright
  \item \emph{Say it:} \emph{podnimát}
  \item \emph{Say it:} \emph{podnimát}, once more
  \item \emph{Say it:} say \emph{podnimát}
  \item \emph{Recall:} say the Russian for to remember, then the Russian for to pull, then say \emph{podnimát} again
\end{itemize}

\begin{tcolorbox}[breakable,colback=teal!4,colframe=teal!35!black,title={Before you move on}]
What does \ru{поднимать} mean? (\textquotedblleft{}To lift, to raise\textquotedblright{}.) Say it once more.
\end{tcolorbox}
\section[\texorpdfstring{\ru{включать} (vklyuchát)}{включать (vklyuchát)}]{\ru{включать} (vklyuchát) — to switch on}

\label{lesson:RU-C175-vklyuchat}



\begin{quote}
Before the new one: say the Russian for to pull, then the Russian for to lift.
\end{quote}

\subsection*{You'll want to know: \ru{включать}}
\textbf{\ru{включать}} (\emph{vklyuchát}) — \textquotedblleft{}to switch on\textquotedblright{}.

\begin{grammarlens}[title={What you've built}]
One more word for this chapter.
\end{grammarlens}

\subsection*{Your turn}
\begin{itemize}
\raggedright
  \item \emph{Say it:} \emph{vklyuchát}
  \item \emph{Say it:} \emph{vklyuchát}, once more
  \item \emph{Say it:} say \emph{vklyuchát}
  \item \emph{Recall:} say the Russian for to pull, then the Russian for to lift, then say \emph{vklyuchát} again
\end{itemize}

\begin{tcolorbox}[breakable,colback=teal!4,colframe=teal!35!black,title={Before you move on}]
What does \ru{включать} mean? (\textquotedblleft{}To switch on\textquotedblright{}.) Say it once more.
\end{tcolorbox}
\section[\texorpdfstring{\ru{выключать} (vyklyuchát)}{выключать (vyklyuchát)}]{\ru{выключать} (vyklyuchát) — to switch off}

\label{lesson:RU-C175-vyklyuchat}



\begin{quote}
Before the new one: say the Russian for to lift, then the Russian for to switch on.
\end{quote}

\subsection*{You'll want to know: \ru{выключать}}
\textbf{\ru{выключать}} (\emph{vyklyuchát}) — \textquotedblleft{}to switch off\textquotedblright{}.

\begin{grammarlens}[title={What you've built}]
One more word for this chapter.
\end{grammarlens}

\subsection*{Your turn}
\begin{itemize}
\raggedright
  \item \emph{Say it:} \emph{vyklyuchát}
  \item \emph{Say it:} \emph{vyklyuchát}, once more
  \item \emph{Say it:} say \emph{vyklyuchát}
  \item \emph{Recall:} say the Russian for to lift, then the Russian for to switch on, then say \emph{vyklyuchát} again
\end{itemize}

\begin{tcolorbox}[breakable,colback=teal!4,colframe=teal!35!black,title={Before you move on}]
What does \ru{выключать} mean? (\textquotedblleft{}To switch off\textquotedblright{}.) Say it once more.
\end{tcolorbox}
\section[\texorpdfstring{\ru{зажигать} (zazhigát)}{зажигать (zazhigát)}]{\ru{зажигать} (zazhigát) — to light, to set alight}

\label{lesson:RU-C175-zazhigat}



\begin{quote}
Before the new one: say the Russian for to switch on, then the Russian for to switch off.
\end{quote}

\subsection*{You'll want to know: \ru{зажигать}}
\textbf{\ru{зажигать}} (\emph{zazhigát}) — \textquotedblleft{}to light, to set alight\textquotedblright{}.

\begin{grammarlens}[title={What you've built}]
One more word for this chapter.
\end{grammarlens}

\subsection*{Your turn}
\begin{itemize}
\raggedright
  \item \emph{Say it:} \emph{zazhigát}
  \item \emph{Say it:} \emph{zazhigát}, once more
  \item \emph{Say it:} say \emph{zazhigát}
  \item \emph{Recall:} say the Russian for to switch on, then the Russian for to switch off, then say \emph{zazhigát} again
\end{itemize}

\begin{tcolorbox}[breakable,colback=teal!4,colframe=teal!35!black,title={Before you move on}]
What does \ru{зажигать} mean? (\textquotedblleft{}To light, to set alight\textquotedblright{}.) Say it once more.
\end{tcolorbox}
